%! AP Statistics — Unit 4, Lesson 4.3 — Student Packet Cover Sheet
\documentclass[10pt]{article}
\usepackage{apstats-article}
\usepackage{apstats-boxes}
\usepackage{ltablex}
\keepXColumns

\begin{document}

% Full-bleed navy banner
\begin{tikzpicture}[remember picture, overlay]
  \fill[navy] (current page.north west)
    rectangle ([yshift=-0.9in]current page.north east);
\end{tikzpicture}
\vspace{-0.8in}

\begin{center}
  {\color{white}\textbf{\LARGE AP Statistics}}\\[4pt]
  {\color{skymid}\large\textbf{Unit 4: Probability, Random Variables, and Probability Distributions}}\\[6pt]
  {\color{white}\textbf{\large Lesson 4.3 \quad Introduction to Probability}}
\end{center}

\vspace{0.22in}
\namedateperiod
\vspace{0.18in}

\begin{learningtargetbox}
\small
\begin{enumerate}[leftmargin=1.5em, itemsep=5pt]
  \item \ldots{} calculate the probability of an event using sample spaces and the classical definition.
  \item \ldots{} apply the complement rule to find $P(E^c) = 1 - P(E)$.
  \item \ldots{} interpret a probability as the long-run relative frequency of an event.
\end{enumerate}
\end{learningtargetbox}

\vspace{0.14in}

\begin{tocbox}
\small
\renewcommand{\arraystretch}{1.5}
\begin{tabularx}{\linewidth}{@{} c l X r @{}}
\rowcolor{sky}
\textbf{\#} & \textbf{Component} & \textbf{Description} & \textbf{Score} \\
\midrule
1 & Warm-Up       & Spiral review + probability preview                    & \blank{1.2cm} \\
2 & Guided Notes  & Sample spaces; classical probability; complement rule   & \blank{1.2cm} \\
3 & Group Activity & Cards/spinner: probability and complement calculations  & \blank{1.2cm} \\
4 & Exit Ticket   & Spinner with 8 equal sections                          & \blank{1.2cm} \\
5 & Homework      & Multi-scenario probability practice                     & \blank{1.2cm} \\
\midrule
  & \multicolumn{2}{r}{\textbf{Total}} & \blank{1.2cm} \\
\end{tabularx}
\end{tocbox}

\vspace{0.14in}

\begin{remindbox}
\small
The \textbf{most important ideas in Lesson 4.3}: probability is a number between 0 and 1
that measures how likely an event is in the long run.

\vspace{6pt}
\begin{tabularx}{\linewidth}{@{} >{\centering\arraybackslash\columncolor{goldbg}}p{0.06\linewidth}
                                    X @{}}
\textbf{\large!} &
\textbf{Classical probability formula.}\enspace
When all outcomes are equally likely:
$P(E) = \dfrac{\text{number of outcomes in }E}{\text{total number of outcomes in }S}$. \\[6pt]
\textbf{\large!} &
\textbf{Complement rule.}\enspace
$P(E^c) = 1 - P(E)$. Often it is easier to calculate the probability of the complement
and subtract from 1. \\[6pt]
\textbf{\large!} &
\textbf{Long-run interpretation.}\enspace
$P(E) = 0.3$ means: if we repeat this random process many times, event $E$ will occur
in about 30\% of the repetitions. \\
\end{tabularx}
\end{remindbox}

\end{document}
